\subsection{A three-element first corner and two small outer stalks}

\begin{theorem}[Proved: two small stalks at a three-element first corner]
\label{CP2:A-linear-three-initial-two-small-stalks-nineteen}
Let $C$ have full linear support and at most nineteen arrows.
Suppose its first corner has order at most three, its middle
corner is trivial, its last corner has order at most eight,
and $|Y|=|C(1,2)|=2$.  Suppose further that
\[
 (|X|,|Z|)=(2,3)\quad\hbox{or}\quad(3,2),
 \qquad X=C(0,1),\quad Z=C(0,2).
\]
Then $BC$ cannot have projective-plane type.
\end{theorem}
\begin{proof}
Assume it has that type.  The complete bound through eighteen
arrows lets us assume there are exactly nineteen.  If the trivial left rational
module of the first corner $A$ is projective, take the
actual exact-orbit quotient at that corner.  This is a
rational ring equivalence, preserves boundedness, leaves
$Y$ unchanged and does not increase $|X|,|Z|$.  Its first
two corners are trivial, so
Proposition~\ref{CP2:A-two-small-incoming-eleven} is a
contradiction.  The elementary order-three classification
leaves only the identity-adjoined two-element left-zero
band $A=\{1,e_0,e_1\}$, with $e_ie_j=e_i$.

Let $P_i$ and $Q_i$ be its right actors on $X$ and $Z$.
Each pair consists of idempotent retractions onto one
common image.  Their compatibility with composition is
\[
 Q_i\mu(y,x)=\mu(y,P_ix).
\]
All $Q_i$ commute with the last-corner action on $Z$.

First suppose $|X|=2$, $|Z|=3$.  Then $P_0=P_1=P$:
two retractions of a two-set onto the same image coincide.
If $Q_0=Q_1=Q$, the maps $P$ on $X$ and $Q$ on $Z$
form an idempotent natural endomorphism of the cut from
object $0$ to objects $1,2$.  If either map is proper,
Lemma~\ref{CP2:A-cut-retract} gives an actual smaller
projective-plane category, contradicting the complete
bound through eighteen arrows.  If both are identities,
the nontrivial first corner acts trivially on every
outgoing stalk.  Its actual trivial-corner retract has
the same rational $H^2$; the integral retract argument
of Corollary~\ref{CP2:A-three-object-trivial-extreme-actions-nineteen}
again gives a smaller projective-plane category.

Suppose $Q_0\ne Q_1$.  Their common image $U\subset Z$
has two elements and $Z\setminus U$ has one member $w$;
the values $Q_0w,Q_1w$ are distinct.  The two maps agree
exactly on $U$, where both are the identity.  Since
$P_0=P_1$, every composite $\mu(y,x)$ lies in $U$.
If $P$ is constant, the compatibility equation makes
$\mu(y,x)=\mu(y,Px)$ for every $x$.  Therefore $P$ on
$X$, together with the identity on $Z$, is a proper
natural cut retraction.  The same eighteen-arrow
contradiction applies.

It remains that $P$ is the identity.  Put
$D=C|_{\{0,1\}}$.  The space $BA$ is contractible and
the action on $X$ is trivial, so the actual edge-bar
pushout identifies $BD$ rationally with a graph.
In particular its rational homology in degrees at
least two is zero.

Use the connecting element category with an arrow
$z\to z'$ labelled $a\in A$ exactly when $z'a=z$,
and an arrow $z\to y$ labelled $x\in X$ exactly
when $\mu(y,x)=z$.  Its objects are $U\cup\{w\}\cup Y$.
Each core object $u_i=Q_iw$ has corner $A$; all other
corners are trivial.  Its strict arrows are one arrow
$u_i\to w$ labelled $e_i$ for each $i$, and the arrows
$\mu(y,x)\to y$ labelled $x$.  There are no other
strict arrows and no composable pair of strict arrows.
The right $A$ action on each of these outgoing stalks
is trivial: $e_i a=e_i$, and $xa=x$.

Collapse the two core copies of $A$ to trivial corners,
keeping every strict arrow.  This is an actual functor
to a category whose nerve is a graph $\Gamma$ with
the two directed legs $u_i\to w$ and the $|Y||X|$
edges $\mu(y,x)\to y$.  Its corner and edge-bar maps
are rational equivalences because $BA$ is contractible
and all these outgoing actions are trivial.  The
height-one bar pushout proves the actual rational
homology equivalence on the full connecting category.
It is natural for the commuting $N$ action.  Indeed,
if $nw=w$, then $nu_i=u_i$; if $nw\in U$, then
$nu_0=nu_1=nw$.  Thus a leg is sent either to itself
or to a vertex.  The other edges keep their $x$
labels and send $y$ to $ny$.  All these rules are
literal functor rules and compose strictly.  They
give exactly the directed-graph action allowed in
Lemma~\ref{CP2:A-directed-graph-collapses-eight}.

Collapsing the actual subspace $BD$ of $BC$ gives the
homotopy cofiber of the connecting homotopy-orbit space
mapping to $BN$.  Its pullback onto $H^2(BC;\mathbb Q)$
is surjective, and it is an isomorphism in degree four,
because $H_j(BD;\mathbb Q)=0$ for $j\geq2$.
The cofiber has bounded rational homology.  Its actual
graph comparison and the cited order-eight theorem
kill all second products.  Their pullbacks contradict
the projective-plane square.

Now suppose $|X|=3$, $|Z|=2$.  This time
$Q_0=Q_1=Q$.  If also $P_0=P_1=P$, the same pair of
maps $P,Q$ is a proper cut retraction unless both are
identities; that last case is the same trivial
extreme-corner retract.  Both are already excluded.

Otherwise the common image of the $P_i$ is a two-set
$\{u_0,u_1\}$, and there is one remaining state $v$
with $P_0v=u_0$, $P_1v=u_1$, after renaming.
If $Q$ is the identity, compatibility gives
\[
 \mu(y,u_0)=\mu(y,v)=\mu(y,u_1)
 \quad\text{for every }y.
\]
The constant map $X\to\{u_0\}$ and the identity on
$Z$ are therefore a proper natural cut retraction.
If $Q$ is constant with value $z_0$, compatibility
instead gives
$\mu(y,u_0)=\mu(y,u_1)=z_0$.
The constant map on $X$ and $Q$ on $Z$ are then a
proper natural cut retraction.  The latter map is
$N$-equivariant because it is a right $A$ actor,
and both maps commute with the right $A$ action.
Each image has fewer than nineteen arrows and would
retain actual projective-plane type.  This finishes
all cases.
\end{proof}

\begin{corollary}[Proved: two further order-eight endpoint profiles]
\label{CP2:A-linear19-three-initial-large-profiles}
The profiles $(3,1,8;2,2,3)$ and $(3,1,8;3,2,2)$,
original indices $441,442$ in the frozen linear
inventory, cannot have projective-plane type.
\end{corollary}
\begin{proof}
Apply the theorem.  Both categories have nineteen arrows.
\end{proof}
